\subsection{平坦性判别准则的证明}

\textbf{(A) 蕴含关系 (i) \(\Rightarrow\) (ii) \(\Leftrightarrow\) (iii)}\par

蕴含关系 (i) \(\Rightarrow\) (ii) 是立即的（第 I 章 § 4）。等价关系 (ii) \(\Leftrightarrow\) (iii) 是第 I 章 § 4 命题 2 的特例；该命题应用于 R = A、S = A/3、F = M、E = N，并注意到：给 N 赋予一个 (A/3)-模结构，等价于给 N 赋予一个 A-模结构，使 N 被 3 消灭。

注 (1)。条件 (ii) 也等价于下述条件：

(ii') \(\operatorname{Tor}_1^{\mathbf{A}}(\mathbf{N},\mathbf{M}) = 0\)，对于每个被 \(\Im\) 的某个幂消灭的 A-模 N 都成立。显然 (ii') 蕴含 (ii)。反之，若 (ii) 成立，则特别有 \(\operatorname{Tor}_1^{\mathbf{A}}(\Im^n\mathrm{N} / \Im^{n + 1}\mathrm{N},\mathbf{M}) = 0\) 对所有 \(n\) 成立；由正合列

\[
0 \rightarrow \Im^ {n + 1} N \rightarrow \Im^ {n} N \rightarrow \Im^ {n} N / \Im^ {n + 1} N \rightarrow 0
\]

得到正合列

\[
\operatorname{Tor} _ {1} ^ {\mathbf {A}} \left(\Im^ {n + 1} \mathrm{N}, \mathrm{M}\right)\rightarrow \operatorname{Tor} _ {1} ^ {\mathbf {A}} \left(\Im^ {n} \mathrm{N}, \mathrm{M}\right)\rightarrow \operatorname{Tor} _ {1} ^ {\mathbf {A}} \left(\Im^ {n} \mathrm{N} / \Im^ {n + 1} \mathrm{N}, \mathrm{M}\right)
\]

又由于存在整数 \(m\) 使得 \(\mathfrak{S}^m\mathrm{N} = 0\)，我们对 \(n\) 作递降归纳，得到 \(\operatorname{Tor}_1^{\mathbf{A}}(\mathfrak{S}^n\mathrm{N},\mathbf{M}) = 0\) 对所有 \(n\leqslant m\) 成立，特别地对 \(n = 0\) 成立。

由此可知，若 \(\Im\) 是幂零的，则 (ii) 蕴含 (i)，因为此时 (ii') 意味着对每个 A-模都有 \(\operatorname{Tor}_1^{\mathbf{A}}(\mathbf{N},\mathbf{M}) = 0\)，其中 \(\mathbf{N}\) 是任意 A-模，从而 \(\mathbf{M}\) 是平坦的（第 I 章 § 4）。

\textbf{(B) 证明下述命题：}\par

命题 1. 设 A 是交换环，\(\mathfrak{I}\) 是 A 的理想，M 是 A-模。以下条件等价：

\begin{enumerate}
\def\labelenumi{(\alph{enumi})}
\item
  对所有 \(n \geqslant 1\)，\(\operatorname{Tor}_1^{\mathbf{A}}(\mathbf{A}/\mathfrak{I}^n, \mathbf{M}) = 0\)。
\item
  对所有 \(n \geqslant 1\)，典范同态
\end{enumerate}

\[
\theta_ {n} \colon \mathfrak {I} ^ {n} \otimes_ {\mathbf {A}} \mathbf {M} \rightarrow \mathfrak {I} ^ {n} \mathbf {M}
\]

是双射。

此外，这些条件蕴含：

\begin{enumerate}
\def\labelenumi{(\alph{enumi})}
\setcounter{enumi}{2}
\tightlist
\item
  典范同态 \(\gamma_{\mathrm{M}}: \operatorname{gr}(\mathbf{A}) \otimes_{\mathbf{g}_{0}(\mathbf{A})} \operatorname{gr}_{0}(\mathbf{M}) \to \operatorname{gr}(\mathbf{M})\) 是双射。反之，若 \(\Im\) 是幂零的，则 (c) 蕴含 (a) 和 (b)。
\end{enumerate}

由正合列可得 (a) 与 (b) 等价：

\[
0 = \operatorname{Tor} _ {1} ^ {\mathbf {A}} (\mathrm{A}, \mathrm{M}) \rightarrow \operatorname{Tor} _ {1} ^ {\mathbf {A}} (\mathrm{A} / \Im^ {n}, \mathrm{M}) \rightarrow \Im^ {n} \otimes_ {\mathbf {A}} \mathrm{M} \rightarrow \mathrm{M}.
\]

再考虑图

\[
\begin{array}{c} \mathfrak {I} ^ {n + 1} \otimes_ {\mathrm{A}} \mathrm{M} \longrightarrow \mathfrak {I} ^ {n} \otimes_ {\mathrm{A}} \mathrm{M} \longrightarrow (\mathfrak {I} ^ {n} / \mathfrak {I} ^ {n + 1}) \otimes_ {\mathrm{A}} (\mathrm{M} / \mathfrak {I M}) \longrightarrow 0 \\ \theta_ {n + 1} \Bigg \downarrow \quad \theta_ {n} \Bigg \downarrow \quad \gamma_ {n} \Bigg \downarrow \\ 0 \longrightarrow \mathfrak {I} ^ {n + 1} \mathrm{M} \longrightarrow \mathfrak {I} ^ {n} \mathrm{M} \longrightarrow \operatorname{gr} _ {n} (\mathrm{M}) \longrightarrow 0 \end{array}\tag{1}
\]

这里注意到，\((\mathfrak{I}^n / \mathfrak{I}^{n+1}) \otimes_{\mathbf{A}} (\mathbf{M} / \mathfrak{I}\mathbf{M})\) 典范地识别为 \((\mathfrak{I}^n / \mathfrak{I}^{n+1}) \otimes_{\mathbf{A} / \mathfrak{S}} (\mathbf{M} / \mathfrak{I}\mathbf{M})\)。根据 \(\gamma_n\) 的定义，此图交换且其各行正合。若 (b) 成立，则 \(\theta_n\) 和 \(\theta_{n+1}\) 是双射，因而由余核的定义 \(\gamma_n\) 也是双射，于是 (b) 蕴含 (c)。反之，假设 \(\mathfrak{I}\) 幂零，我们证明 (c) 蕴含 (b)；我们对 \(n\) 作递降归纳，因为 \(\mathfrak{I}^n \otimes_{\mathbf{A}} \mathbf{M} = \mathfrak{I}^n\mathbf{M} = 0\) 在 \(n\) 足够大时成立。于是，若在图 (1) 中 \(\gamma_n\) 和 \(\theta_{n+1}\) 是双射，则根据第 I 章 § 1 第 4 节命题 2 的推论 1，\(\theta_n\) 也是双射。

\textbf{(C) 蕴含关系 (ii) \(\Rightarrow\) (iv)}\par

若 (ii) 成立，则根据注 1，(ii') 也成立；于是命题 1 表明 \(\gamma_{\mathbf{M}}\) 是同构。另一方面，我们已经知道 (ii) 蕴含 (iii)，因而 M/Im 是平坦的 (A/Im)-模，这就完成了 (ii) 蕴含 (iv) 的证明。

注 (2)。命题 1 表明，若 \(\mathfrak{S}\) 是幂零的，则 (iv) 蕴含 (iii)；结合注 1，我们因而证明了在此情形下 (i)、(ii)、(iii) 和 (iv) 等价。

\textbf{(D) 等价关系 (iv) \(\Leftrightarrow\) (v)}\par

  对所有 \(n \geqslant 1\)，M/ \(\Im^{n}M\) 具有典范的 (A/ \(\Im^{n}\) )-模结构。若以 \((\Im/\Im^{n})\)-进滤过对其进行滤过，立即有：\(\mathrm{gr}_{m}(\mathrm{M}/\Im^{n}\mathrm{M}) = \mathrm{gr}_{m}(\mathrm{M})\) 当 m \textless{} n，且 \(\mathrm{gr}_{m}(\mathrm{M}/\Im^{n}\mathrm{M}) = 0\) 当 \(m \geqslant n\)。对于所有 \(k \geqslant 1\)，令 \(A_{k} = A/\Im^{k}\)、\(\Im_{k} = \Im/\Im^{k} M_{k} = M/\Im^{k} M\)；令 \((\mathrm{iv})_{k}\)（分别为 \((\mathrm{v})_{k}\)）表示将 A、\(\Im\)、M 替换为 \(A_{k}\)、\(\Im_{k}\)、\(M_{k}\) 后由 (iv)（分别由 (v)）得到的断言。由刚才所述，(iv) 等价于“对所有 \(k \geqslant 1\)，\((\mathrm{iv})_{k}\)”；显然 (v) 等价于“对所有 \(k \geqslant 1\)，\((\mathrm{v})_{k}\)”。于是只需对所有 k 证明 \((\mathrm{iv})_{k} \Leftrightarrow (\mathrm{v})_{k}\)，或者只需证明 \(\Leftrightarrow (\mathrm{v})\) 当 \(\Im\) 幂零时。现在（注 2）我们已经看到，在此情形下 (iv) 等价于 (i)。由于 M/ \(\Im^{n}M\) 同构于 M \(\otimes_{A} (A/\Im^{n})\)，(i) 蕴含 (v)（第 I 章 §2，第 7 节，命题 8 的推论 2）；此外显然 (v) 蕴含 (i)。因此，我们已在所有情形证明了 (iv) \(\Leftrightarrow (\mathrm{v})\)，并且在 \(\Im\) 幂零时证明了定理的所有性质都等价。

\textbf{(E) 当 \((\mathbf{v})\Rightarrow (\mathbf{i})\) 且 \(\mathbf{A}\) 是 Noether 环、\(\mathbf{M}\) 是理想 Hausdorff 模时}\par

只需证明，对于每个理想 \(\mathfrak{a}\)（它属于 \(\mathbf{A}\)），典范映射 \(j\colon \mathfrak{a} \otimes_{\mathbf{A}} \mathbf{M} \to \mathbf{M}\) 是单射（第 I 章 § 2，第 3 节，命题 1）。令 \(x \in \operatorname{Ker} j\)；由于 \(\mathfrak{a} \otimes_{\mathbf{A}} \mathbf{M}\) 关于 \(\mathfrak{S}\)-进拓扑是 Hausdorff 的，只需验证对于每个整数 \(n > 0\)，都有 \(x \in \mathfrak{S}^n(\mathfrak{a} \otimes_{\mathbf{A}} \mathbf{M})\)。令 \(f\colon \mathfrak{S}^n\mathfrak{a} \to \mathfrak{a}\) 为典范嵌入；只需证明 \(x \in \operatorname{Im}(f \otimes 1_{\mathbf{M}})\)；因为若 \(b \in \mathfrak{S}^n\)、\(a \in \mathfrak{a}\) 且 \(m \in \mathbf{M}\)，则 \(f \otimes 1_{\mathbf{M}}\) 作用于元素 \((ba) \otimes m\)（该元素属于 \((\mathfrak{S}^n\mathfrak{a}) \otimes_{\mathbf{A}} \mathbf{M}\)）所得的像，是元素 \((ba) \otimes m = b(a \otimes m)\)，它属于 \(\mathfrak{a} \otimes_{\mathbf{A}} \mathbf{M}\)，故 \(\operatorname{Im}(f \otimes 1_{\mathbf{M}}) \subset \mathfrak{S}^n(\mathfrak{a} \otimes_{\mathbf{A}} \mathbf{M})\)。由 Krull 定理（§ 3，第 2 节，定理 2），存在整数 \(k\)，使得 \(\mathfrak{a}_k = \mathfrak{a} \cap \mathfrak{S}^k \subset \mathfrak{S}^n\mathfrak{a}\)；若 \(i: \mathfrak{a}_k \to \mathfrak{a}\) 是典范嵌入，则只需证明 \(x \in \operatorname{Im}(i \otimes 1_M)\)。现在，记 \(p: \mathfrak{a} \to \mathfrak{a}/\mathfrak{a}_k\) 和 \(h: \mathfrak{a}/\mathfrak{a}_k \to \mathrm{A}/\mathfrak{S}^k\) 为典范映射，则存在交换图

\[
\begin{array}{c} \mathfrak {a} _ {k} \otimes_ {\mathrm{A}} \mathrm{M} \xrightarrow {i \otimes 1 _ {\mathrm{M}}} \mathfrak {a} \otimes_ {\mathrm{A}} \mathrm{M} \xrightarrow {p \otimes 1 _ {\mathrm{M}}} (\mathfrak {a} / \mathfrak {a} _ {k}) \otimes_ {\mathrm{A}} \mathrm{M} \longrightarrow 0 \\ j \Biggl \downarrow \\ \mathrm{M} \quad \longrightarrow \quad (\mathrm{A} / \mathfrak {I} ^ {k}) \otimes_ {\mathrm{A}} \mathrm{M} \end{array}
\]

其中第一行正合。只需证明 \(x \in \operatorname{Ker}(p \otimes 1_{\mathbf{M}})\)；又由于按假设 \(x \in \operatorname{Ker} j\)，故只需验证映射 \(h \otimes 1_{\mathbf{M}}\) 是单射。现在，它还可写成（《代数》第 II 章 § 3 第 6 节命题 6 的推论 3）

\[
h \otimes 1 _ {\mathbf {M} / \Im^ {k} \mathbf {M}}: (a / a _ {k}) \otimes_ {\mathbf {A} / \Im^ {k}} (\mathbf {M} / \Im^ {k} \mathbf {M}) \rightarrow \mathbf {M} / \Im^ {k} \mathbf {M}
\]

由于 \(h\) 是单射，且根据 (v)，\(\mathbf{M} / \Im^k\mathbf{M}\) 是平坦的 \((\mathbf{A} / \Im^k)\)-模，证明完毕。

